% GENERATED FILE. Edit canonical lessons, then run npm run generate:books.
% canonical-source-hash: fnv1a64:7790144bd6e29163
% canonical-lessons: TE-C278-pattuko, TE-C278-snanam-ceyu, TE-C278-tayaru-ceyu, TE-C278-datu, TE-C278-nadupu

\chapter{To hold, To bathe, To prepare, To cross, To drive}
\label{ch:te-to-hold-to-bathe-to-prepare-to-cross-to-drive}
\hlchaptermodality{278}

\begin{chapteropening}
\textbf{By the end of this chapter:} \emph{I can say to hold, to bathe, to prepare, to cross and to drive.}

Complete the last lesson of chapter 278: I can say to hold, to bathe, to prepare, to cross and to drive.
\end{chapteropening}

\section[\texorpdfstring{\te{పట్టుకో} (paṭṭukō)}{పట్టుకో (paṭṭukō)}]{\te{పట్టుకో} (paṭṭukō) — to hold, to catch}

\label{lesson:TE-C278-pattuko}



\begin{quote}
Before the new one: say the Telugu for a potter, then the Telugu for a blacksmith.
\end{quote}

\subsection*{You'll want to know: \te{పట్టుకో}}
\textbf{\te{పట్టుకో}} — \emph{paṭṭukō} — \textquotedblleft{}to hold, to catch\textquotedblright{}.

\begin{grammarlens}[title={What you've built}]
One more word for this chapter.
\end{grammarlens}

\subsection*{Your turn}
\begin{itemize}
\raggedright
  \item \emph{Say it:} \emph{paṭṭukō}
  \item \emph{Say it:} \emph{paṭṭukō}, once more
  \item \emph{Say it:} say \emph{paṭṭukō}
  \item \emph{Recall:} say the Telugu for a potter, then the Telugu for a blacksmith, then say \emph{paṭṭukō} again
\end{itemize}

\begin{tcolorbox}[breakable,colback=teal!4,colframe=teal!35!black,title={Before you move on}]
What does \te{పట్టుకో} mean? (\textquotedblleft{}To hold, to catch\textquotedblright{}.) Say it once more.
\end{tcolorbox}
\section[\texorpdfstring{\te{స్నానం} \te{చేయు} (snānaṁ cēyu)}{స్నానం చేయు (snānaṁ cēyu)}]{\te{స్నానం} \te{చేయు} (snānaṁ cēyu) — to bathe}

\label{lesson:TE-C278-snanam-ceyu}



\begin{quote}
Before the new one: say the Telugu for a blacksmith, then the Telugu for to hold.
\end{quote}

\subsection*{You'll want to know: \te{స్నానం} \te{చేయు}}
\textbf{\te{స్నానం} \te{చేయు}} — \emph{snānaṁ cēyu} — \textquotedblleft{}to bathe\textquotedblright{}.

\begin{grammarlens}[title={What you've built}]
One more word for this chapter.
\end{grammarlens}

\subsection*{Your turn}
\begin{itemize}
\raggedright
  \item \emph{Say it:} \emph{snānaṁ cēyu}
  \item \emph{Say it:} \emph{snānaṁ cēyu}, once more
  \item \emph{Say it:} say \emph{snānaṁ cēyu}
  \item \emph{Recall:} say the Telugu for a blacksmith, then the Telugu for to hold, then say \emph{snānaṁ cēyu} again
\end{itemize}

\begin{tcolorbox}[breakable,colback=teal!4,colframe=teal!35!black,title={Before you move on}]
What does \te{స్నానం} \te{చేయు} mean? (\textquotedblleft{}To bathe\textquotedblright{}.) Say it once more.
\end{tcolorbox}
\section[\texorpdfstring{\te{తయారు} \te{చేయు} (tayāru cēyu)}{తయారు చేయు (tayāru cēyu)}]{\te{తయారు} \te{చేయు} (tayāru cēyu) — to prepare, to make}

\label{lesson:TE-C278-tayaru-ceyu}



\begin{quote}
Before the new one: say the Telugu for to hold, then the Telugu for to bathe.
\end{quote}

\subsection*{You'll want to know: \te{తయారు} \te{చేయు}}
\textbf{\te{తయారు} \te{చేయు}} — \emph{tayāru cēyu} — \textquotedblleft{}to prepare, to make\textquotedblright{}.

\begin{grammarlens}[title={What you've built}]
One more word for this chapter.
\end{grammarlens}

\subsection*{Your turn}
\begin{itemize}
\raggedright
  \item \emph{Say it:} \emph{tayāru cēyu}
  \item \emph{Say it:} \emph{tayāru cēyu}, once more
  \item \emph{Say it:} say \emph{tayāru cēyu}
  \item \emph{Recall:} say the Telugu for to hold, then the Telugu for to bathe, then say \emph{tayāru cēyu} again
\end{itemize}

\begin{tcolorbox}[breakable,colback=teal!4,colframe=teal!35!black,title={Before you move on}]
What does \te{తయారు} \te{చేయు} mean? (\textquotedblleft{}To prepare, to make\textquotedblright{}.) Say it once more.
\end{tcolorbox}
\section[\texorpdfstring{\te{దాటు} (dāṭu)}{దాటు (dāṭu)}]{\te{దాటు} (dāṭu) — to cross}

\label{lesson:TE-C278-datu}



\begin{quote}
Before the new one: say the Telugu for to bathe, then the Telugu for to prepare.
\end{quote}

\subsection*{You'll want to know: \te{దాటు}}
\textbf{\te{దాటు}} — \emph{dāṭu} — \textquotedblleft{}to cross\textquotedblright{}.

\begin{grammarlens}[title={What you've built}]
One more word for this chapter.
\end{grammarlens}

\subsection*{Your turn}
\begin{itemize}
\raggedright
  \item \emph{Say it:} \emph{dāṭu}
  \item \emph{Say it:} \emph{dāṭu}, once more
  \item \emph{Say it:} say \emph{dāṭu}
  \item \emph{Recall:} say the Telugu for to bathe, then the Telugu for to prepare, then say \emph{dāṭu} again
\end{itemize}

\begin{tcolorbox}[breakable,colback=teal!4,colframe=teal!35!black,title={Before you move on}]
What does \te{దాటు} mean? (\textquotedblleft{}To cross\textquotedblright{}.) Say it once more.
\end{tcolorbox}
\section[\texorpdfstring{\te{నడుపు} (naḍupu)}{నడుపు (naḍupu)}]{\te{నడుపు} (naḍupu) — to drive, to ride (a vehicle)}

\label{lesson:TE-C278-nadupu}



\begin{quote}
Before the new one: say the Telugu for to prepare, then the Telugu for to cross.
\end{quote}

\subsection*{You'll want to know: \te{నడుపు}}
\textbf{\te{నడుపు}} — \emph{naḍupu} — \textquotedblleft{}to drive, to ride (a vehicle)\textquotedblright{}.

\begin{grammarlens}[title={What you've built}]
One more word for this chapter.
\end{grammarlens}

\subsection*{Your turn}
\begin{itemize}
\raggedright
  \item \emph{Say it:} \emph{naḍupu}
  \item \emph{Say it:} \emph{naḍupu}, once more
  \item \emph{Say it:} say \emph{naḍupu}
  \item \emph{Recall:} say the Telugu for to prepare, then the Telugu for to cross, then say \emph{naḍupu} again
\end{itemize}

\begin{tcolorbox}[breakable,colback=teal!4,colframe=teal!35!black,title={Before you move on}]
What does \te{నడుపు} mean? (\textquotedblleft{}To drive, to ride (a vehicle)\textquotedblright{}.) Say it once more.
\end{tcolorbox}
